\documentclass[12pt]{amsart}
\pagestyle{plain}

\usepackage{amsmath,amssymb}
\usepackage{enumerate}
\usepackage[shortlabels]{enumitem}
\setlist[enumerate]{itemsep=6pt, topsep=6pt, parsep=0pt}
\usepackage{mathtools}
\usepackage{xcolor}
\usepackage[left=0.75in,right=1in,bottom=0.9in]{geometry}
\usepackage[T1]{fontenc}
\usepackage[parfill]{parskip}
\renewcommand{\baselinestretch}{1}

\newcommand{\norm}[1]{\left\|#1\right\|}

% SECTION DIVIDERS
\newcommand{\divider}{
  \vspace{-1em}
  \noindent
  \raisebox{-0.15ex}{\textbullet}%
  \hspace{0.5em}%
  \rule[0.5ex]{0.95\textwidth}{1.0pt}%
  \hspace{0.5em}%
  \raisebox{-0.15ex}{\textbullet}%
  \vspace{0em}
}

\title{ESE 2030 QUIZZAM 3 : SOLUTIONS GUIDE \\ Spring 2026}
\author{prof-g}

\begin{document}
\maketitle

\begin{center}
{\bf NOTE TO STUDENTS}
\end{center}

This solutions guide is organized by course week and topic.
Since there are five versions of this quiz with permuted problem orders \emph{and}
permuted answer choices, the guide does \emph{not} identify correct answers by
letter~(A)--(E).
Instead, the correct answer is \textcolor{blue}{\textbf{highlighted in blue}}
directly within each list of choices.
Match the problem text to the one you received, then read the highlighted
choice, the explanation, and the analysis of each distractor.

Each solution includes:
\begin{itemize}
\item The correct answer (highlighted in blue in the choice list)
\item Detailed mathematical explanation
\item Analysis of every wrong choice, including partial credit designations and common misconceptions
\end{itemize}

\newpage

%%%%%%%%%%%%%%%%%%%%%%%%%%%%%%%%%%%%%%%%%%%%%%%%%%%%%%%%%%%
\section*{WEEK 7: DIAGONALIZATION \& DYNAMICS}
%%%%%%%%%%%%%%%%%%%%%%%%%%%%%%%%%%%%%%%%%%%%%%%%%%%%%%%%%%%

%----------------------------------------------------------
\subsection*{Problem: Superposition of solutions to $d\mathbf{x}/dt = A\mathbf{x}$}
%----------------------------------------------------------

Consider the system $\dfrac{d\mathbf{x}}{dt} = A\mathbf{x}$ where $A$ is a
$3 \times 3$ matrix. Suppose you know two solutions:
\begin{itemize}
\item $\mathbf{x}_1(t)$ is the solution with initial condition $\mathbf{x}_1(0) = \mathbf{a}$
\item $\mathbf{x}_2(t)$ is the solution with initial condition $\mathbf{x}_2(0) = \mathbf{b}$
\end{itemize}
What is the solution $\mathbf{x}(t)$ with initial condition $\mathbf{x}(0) = 3\mathbf{a} - 2\mathbf{b}$?

\begin{enumerate}[(A)]
\item \textcolor{blue}{\textbf{$3\mathbf{x}_1(t) - 2\mathbf{x}_2(t)$}}
\item $e^{At}(3\mathbf{a} - 2\mathbf{b})$, but this cannot be simplified
      further without computing $e^{At}$
\item $(3-2t)\mathbf{x}_1(t) + t\,\mathbf{x}_2(t)$, by interpolation
      between the two known solutions
\item $3\mathbf{x}_1(t) - 2\mathbf{x}_2(t)$, but only if $A$ is diagonalizable
\item Cannot be determined without knowing $A$ explicitly
\end{enumerate}

\textbf{Explanation:}
The solution operator $\mathbf{x}_0 \mapsto e^{At}\mathbf{x}_0$ is
\emph{linear} in the initial condition, because matrix multiplication
is linear. Therefore:
$$e^{At}(3\mathbf{a} - 2\mathbf{b})
  = 3\,e^{At}\mathbf{a} - 2\,e^{At}\mathbf{b}
  = 3\mathbf{x}_1(t) - 2\mathbf{x}_2(t)$$
No knowledge of $A$ is needed beyond the fact that the system is linear.
This is the superposition principle: any linear combination of solutions
is again a solution, with the corresponding linear combination of initial
conditions. This holds for \emph{any} matrix $A$ --- diagonalizable or
not, real or complex eigenvalues, repeated or distinct.

\textbf{Partial credit:}
The choice stating $e^{At}(3\mathbf{a} - 2\mathbf{b})$ ``cannot be
simplified'' is technically a true expression --- it correctly writes
down the formal solution --- but fails to recognize that linearity of
$e^{At}$ lets us rewrite it using the known solutions. The student
understands the matrix exponential formula but misses the deeper point
that known solutions can be combined without ever computing $e^{At}$
explicitly.

\textbf{Trick answer:}
The choice requiring $A$ to be diagonalizable conflates the techniques
for \emph{computing} $e^{At}$ (which do depend on eigenstructure) with
the structural fact that $e^{At}$ is linear (which requires no such
assumption).

\textbf{Trick answer:}
The ``interpolation'' choice invents time-dependent coefficients
$(3-2t)$ and $t$. A linear combination of solutions with
\emph{time-dependent} coefficients is generally not a solution ---
only constant-coefficient combinations are guaranteed to work.

\textbf{Trick answer:}
The choice ``cannot be determined without knowing $A$'' fails to
recognize that linearity allows us to build new solutions from old ones
without any information about $A$ itself.

\divider

%----------------------------------------------------------
\subsection*{Problem: Matrix exponential of a nilpotent matrix}
%----------------------------------------------------------

Which of the following is $e^{A}$ where
$A = \begin{bmatrix} 0 & -2 & 0 \\ 0 & 0 & -2 \\ 0 & 0 & 0 \end{bmatrix}$?

\noindent
The five answer choices are matrices of the form
$\begin{bmatrix} 1 & -2 & * \\ 0 & 1 & -2 \\ 0 & 0 & 1 \end{bmatrix}$
with various $(1,3)$ entries (and one choice involving $e^{-2}$).

\textcolor{blue}{\textbf{Correct answer:}
$\begin{bmatrix} 1 & -2 & 2 \\ 0 & 1 & -2 \\ 0 & 0 & 1 \end{bmatrix}$}

\textbf{Explanation:}
The matrix $A$ is strictly upper triangular, hence nilpotent.
We compute:
$$A^2 = \begin{bmatrix} 0 & 0 & 4 \\ 0 & 0 & 0 \\ 0 & 0 & 0 \end{bmatrix},
\qquad A^3 = 0$$
Since $A^k = 0$ for all $k \geq 3$, the power series terminates:
$$e^A = I + A + \tfrac{1}{2!}A^2
     = \begin{bmatrix} 1&0&0\\0&1&0\\0&0&1 \end{bmatrix}
     + \begin{bmatrix} 0&-2&0\\0&0&-2\\0&0&0 \end{bmatrix}
     + \tfrac{1}{2}\begin{bmatrix} 0&0&4\\0&0&0\\0&0&0 \end{bmatrix}
     = \begin{bmatrix} 1&-2&2\\0&1&-2\\0&0&1 \end{bmatrix}$$
The problem tests the entire chain: recognizing nilpotency so the series
is finite, computing $A^2$ correctly (including $(-2)(-2) = +4$), and
applying the $\frac{1}{k!}$ factorial coefficients.

\textbf{Trick answer:}
The choice with $(1,3)$ entry $= 4$ correctly computes $A^2$ but
forgets the factorial denominator $\frac{1}{2!}$ in the series,
effectively computing $\sum A^k$ instead of $e^A = \sum \frac{1}{k!}A^k$.

\textbf{Trick answer:}
The choice with $(1,3)$ entry $= 0$ computes only $I + A$, missing
the $A^2$ term entirely. This student either assumes $A^2 = 0$
without checking, or does not realize higher-order terms can contribute.

\textbf{Trick answer:}
The choice with $(1,3)$ entry $= -2$ correctly sets up
$I + A + \frac{1}{2}A^2$ but makes a sign error in computing $A^2$:
gets $(-2)(-2) = -4$ instead of $+4$.

\textbf{Trick answer:}
The choice involving $e^{-2}$ entries applies the scalar exponential
function to individual matrix entries, a fundamental misunderstanding ---
the matrix exponential is defined via the power series $\sum \frac{1}{k!}A^k$,
not by exponentiating entries.

\divider

%----------------------------------------------------------
\subsection*{Problem: Eigenvectors for distinct eigenvalues}
%----------------------------------------------------------

Let $A$ be a $3 \times 3$ matrix with eigenvalues $\lambda_1 = 5$,
$\lambda_2 = -2$, and $\lambda_3 = 5$ (note: $\lambda_1 = \lambda_3$).
Let $\mathbf{v}_1$ be an eigenvector for $\lambda_1 = 5$ and $\mathbf{v}_2$
an eigenvector for $\lambda_2 = -2$.
Which of the following \textbf{must} be true?

\begin{enumerate}[(A)]
\item $\mathbf{v}_1$ and $\mathbf{v}_2$ are orthogonal
\item $\mathbf{v}_1 + \mathbf{v}_2$ is an eigenvector of $A$ with
      eigenvalue $5 + (-2) = 3$
\item \textcolor{blue}{\textbf{$\mathbf{v}_1$ and $\mathbf{v}_2$ are linearly independent}}
\item $A$ must be diagonalizable, since it has three eigenvalues
      counting multiplicity
\item Any nonzero scalar multiple $c\mathbf{v}_1$ is also an eigenvector
      of $A$, but with eigenvalue $5c$
\end{enumerate}

\textbf{Explanation:}
Eigenvectors corresponding to \emph{distinct} eigenvalues are always
linearly independent. Since $\lambda_1 = 5 \neq -2 = \lambda_2$, the
eigenvectors $\mathbf{v}_1$ and $\mathbf{v}_2$ must be linearly
independent, regardless of any other properties of $A$. This is a
fundamental theorem: eigenspaces for distinct eigenvalues intersect
only at $\mathbf{0}$.

\textbf{Trick answer:}
Orthogonality requires symmetry of $A$ (via the Spectral Theorem),
which is not given. Distinct eigenvalues guarantee independence, not
orthogonality.

\textbf{Trick answer:}
Eigenvalues do not ``add'' when eigenvectors are summed.
We have $A(\mathbf{v}_1 + \mathbf{v}_2) = 5\mathbf{v}_1 + (-2)\mathbf{v}_2$,
which is not a scalar multiple of $\mathbf{v}_1 + \mathbf{v}_2$ when
$5 \neq -2$. So $\mathbf{v}_1 + \mathbf{v}_2$ is generally \emph{not}
an eigenvector at all.

\textbf{Trick answer:}
Having three eigenvalues counting multiplicity is necessary but not
sufficient for diagonalizability. The repeated eigenvalue $\lambda = 5$
has algebraic multiplicity~$2$; $A$ is diagonalizable only if the
eigenspace for $\lambda = 5$ is also two-dimensional.

\textbf{Trick answer:}
Scalar multiples of eigenvectors are indeed eigenvectors, but the
eigenvalue does \emph{not} change:
$A(c\mathbf{v}_1) = cA\mathbf{v}_1 = c(5\mathbf{v}_1) = 5(c\mathbf{v}_1)$.
The eigenvalue remains~$5$, not $5c$. The eigenvalue is a property of
the \emph{direction}, not the magnitude.

\divider

%----------------------------------------------------------
\subsection*{Problem: What diagonalizability tells us beyond eigenvalues}
%----------------------------------------------------------

A $3 \times 3$ matrix $A$ has eigenvalues $\lambda_1 = 3$, $\lambda_2 = 3$,
and $\lambda_3 = -1$. {\bf Fact:} $A$ is diagonalizable.
What does diagonalizability of $A$ tell us that the eigenvalues alone do not?

\begin{enumerate}[(A)]
\item That $\operatorname{tr}(A) = 5$
\item That $A$ is symmetric, since only symmetric matrices can have
      repeated eigenvalues and still be diagonalizable
\item \textcolor{blue}{\textbf{That there exists a basis of $\mathbb{R}^3$
      consisting entirely of eigenvectors of $A$}}
\item That $\det(A) = -9$
\item That $A$ is invertible, since $e^A$ must be well-defined
      and nonsingular
\end{enumerate}

\textbf{Explanation:}
The eigenvalues alone tell us $\operatorname{tr}(A) = 3 + 3 + (-1) = 5$
and $\det(A) = 3 \cdot 3 \cdot (-1) = -9$; these hold for any
$3 \times 3$ matrix with these eigenvalues, diagonalizable or not.
Diagonalizability adds the crucial information that the eigenspace for
$\lambda = 3$ is two-dimensional (geometric multiplicity matches
algebraic multiplicity), so there exist two independent eigenvectors
for $\lambda = 3$ together with one for $\lambda = -1$, forming a full
eigenbasis for $\mathbb{R}^3$. This is the defining content of
diagonalizability: $\mathbb{R}^3$ has a basis in which $A$ acts by pure
scaling on each basis vector.

\textbf{Trick answer:}
The trace equals the sum of eigenvalues for \emph{every} square matrix,
not just diagonalizable ones. This is eigenvalue-level information.

\textbf{Trick answer:}
Symmetry is strictly stronger than diagonalizability. A matrix can be
diagonalizable without being symmetric (e.g., any matrix with distinct
eigenvalues). The implication goes only one way:
symmetric $\Rightarrow$ diagonalizable.

\textbf{Trick answer:}
The determinant equals the product of eigenvalues for every square
matrix, regardless of diagonalizability.

\textbf{Trick answer:}
Since no eigenvalue is zero, $\det(A) = -9 \neq 0$, so invertibility
already follows from the eigenvalues alone. Moreover, $e^A$ is always
nonsingular (for any $A$, since
$\det(e^A) = e^{\operatorname{tr}(A)} \neq 0$), so this tells us nothing
new.

\divider

%----------------------------------------------------------
\subsection*{Problem: Fibonacci sequence and dominant eigenvalue}
%----------------------------------------------------------

The Fibonacci sequence $F_0 = 0,\; F_1 = 1,\; F_{k+1} = F_k + F_{k-1}$
can be written in matrix form:
$$\begin{pmatrix} F_{k+1} \\ F_k \end{pmatrix}
  = \underbrace{\begin{bmatrix} 1 & 1 \\ 1 & 0 \end{bmatrix}}_{A}
    \begin{pmatrix} F_k \\ F_{k-1} \end{pmatrix}
\quad \Rightarrow \quad
\lambda_1 = \dfrac{1+\sqrt{5}}{2} \approx 1.618\ldots\,;\;\;
\lambda_2 = \dfrac{1-\sqrt{5}}{2} \approx -0.618\ldots$$
It is a classical fact that $\dfrac{F_{k+1}}{F_k} \to \lambda_1$.
Which of the following best explains \textbf{why}?

\begin{enumerate}[(A)]
\item $\lambda_1$ is the only positive eigenvalue of $A$, and Fibonacci
      numbers are positive, so the ratio must converge to the positive
      eigenvalue
\item The trace of $A$ is $1$, so the sum of eigenvalues is $1$; as
      $k$ grows, this forces the ratio toward the larger eigenvalue
\item $A$ is symmetric, so the Spectral Theorem guarantees that the
      dominant eigenvalue controls the long-term ratio
\item $\lambda_1$ satisfies $\lambda_1^2 = \lambda_1 + 1$, the same
      relation as $F_{k+1} = F_k + F_{k-1}$ after dividing by $F_k$;
      therefore any sequence obeying this recurrence must converge to
      $\lambda_1$
\item \textcolor{blue}{\textbf{Since $|\lambda_2| < 1 < |\lambda_1|$,
      the component of $\mathbf{x}_k = A^k\mathbf{x}_0$ along the
      subdominant eigenvector decays relative to the component along
      $\lambda_1$'s eigenvector, so $\mathbf{x}_k$ becomes
      approximately parallel to $\lambda_1$'s eigenvector, and
      multiplication by $A$ then scales by approximately $\lambda_1$}}
\end{enumerate}

\textbf{Explanation:}
Expanding in the eigenbasis:
$\mathbf{x}_k = c_1\lambda_1^k\mathbf{v}_1 + c_2\lambda_2^k\mathbf{v}_2$.
Since $|\lambda_2| \approx 0.618 < 1$ and $|\lambda_1| \approx 1.618 > 1$,
the first term dominates: for large $k$,
$\mathbf{x}_k \approx c_1\lambda_1^k\mathbf{v}_1$, so both entries of
$\mathbf{x}_k$ grow by factor $\lambda_1$ at each step, making
$F_{k+1}/F_k \to \lambda_1$. This is exactly the mechanism behind the
\emph{power method}: repeated multiplication aligns the iterate with the
dominant eigenvector, and successive ratios converge to the dominant
eigenvalue.

\textbf{Trick answer:}
The algebraic observation that $\lambda_1^2 = \lambda_1 + 1$ correctly
identifies the limit \emph{assuming convergence}, but does not explain
\emph{why} convergence occurs. Without the eigenvalue magnitude analysis,
one cannot rule out oscillation or divergence. This is a necessary condition
masquerading as a sufficient explanation.

\textbf{Trick answer:}
$A$ is indeed symmetric, but symmetry is not the mechanism behind
convergence. Any diagonalizable matrix with a strictly dominant eigenvalue
exhibits the same behavior --- the convergence comes from eigenvalue
magnitude separation, which works identically for non-symmetric matrices.

\textbf{Trick answer:}
The trace constrains the eigenvalue \emph{sum} but provides no mechanism
for ``forcing'' the ratio toward the larger eigenvalue. Convergence is
governed by eigenvalue magnitudes, not their sum.

\textbf{Trick answer:}
Positivity of $\lambda_1$ and positivity of Fibonacci numbers are both
true, but convergence has nothing to do with sign matching. The operative
criterion is the \emph{magnitude} of the eigenvalue, not its sign.

\divider

%----------------------------------------------------------
\subsection*{Problem: Decoupling via diagonalization}
%----------------------------------------------------------

Consider the coupled system $\dfrac{d\mathbf{x}}{dt} = A\mathbf{x}$ where
$A$ is a $2 \times 2$ diagonalizable matrix with distinct real eigenvalues
$\lambda_1 \neq \lambda_2$ and corresponding eigenvectors
$\mathbf{v}_1, \mathbf{v}_2$.
We solve by writing $A = V\Lambda V^{-1}$ and substituting
$\mathbf{x} = V\mathbf{y}$ to obtain
$\dfrac{d\mathbf{y}}{dt} = \Lambda\mathbf{y}$.
What is the essential reason this change of variables helps?

\begin{enumerate}[(A)]
\item It orthogonalizes the eigenvectors, which is necessary for the
      solution components to evolve independently
\item It makes the matrix $\Lambda$ invertible, which is required to
      compute the solution $\mathbf{x}(t) = e^{At}\mathbf{x}(0)$
\item It eliminates the eigenvalues, replacing them with simpler constants
      that are easier to exponentiate
\item \textcolor{blue}{\textbf{Since $\Lambda$ is diagonal, the
      transformed system consists of independent scalar equations
      $dy_i/dt = \lambda_i y_i$, each of which can be solved separately}}
\item It converts the differential equation into an algebraic equation,
      removing the need to compute any exponentials
\end{enumerate}

\textbf{Explanation:}
The fundamental benefit of diagonalization is \emph{decoupling}. In the
original coordinates, each component's derivative depends on the other
components. After substituting $\mathbf{x} = V\mathbf{y}$, the diagonal
structure of $\Lambda$ breaks the system into independent scalar equations
$dy_i/dt = \lambda_i y_i$, each with immediate solution
$y_i(t) = y_i(0)e^{\lambda_i t}$ from one-variable calculus.
Diagonalization transforms a coupled system (hard) into independent
scalar systems (easy).

\textbf{Partial credit:}
The choice about orthogonalizing eigenvectors correctly identifies
\emph{independence of solution components} as the goal, but incorrectly
attributes this to orthogonality. The eigenvectors need not be orthogonal
at all --- it is the \emph{diagonality} of $\Lambda$ that decouples the
equations, not any orthogonality property of $V$.

\textbf{Trick answer:}
$\Lambda$ may or may not be invertible (an eigenvalue could be zero),
and invertibility is irrelevant. The matrix exponential $e^{At}$ is always
well-defined regardless.

\textbf{Trick answer:}
The eigenvalues are not eliminated --- they appear directly as the
diagonal entries of $\Lambda$ and in the solutions $e^{\lambda_i t}$.
The change of variables \emph{isolates} the eigenvalues; it does not
remove them.

\textbf{Trick answer:}
The problem remains a differential equation after the substitution ---
it just becomes a simpler one. We still compute exponentials
($e^{\lambda_i t}$), but they are scalar exponentials rather than
matrix exponentials.

\divider

%----------------------------------------------------------
\subsection*{Problem: Zero eigenvalue and the kernel}
%----------------------------------------------------------

A $4 \times 4$ matrix $A$ has characteristic polynomial
$p(\lambda) = \lambda^2(\lambda - 3)(\lambda + 1)$.
Which of the following must be true?

\begin{enumerate}[(A)]
\item \textcolor{blue}{\textbf{$\ker(A) \neq \{\mathbf{0}\}$}}
\item $A^{-1}$ exists
\item $\operatorname{tr}(A) = 0$
\item $e^{A}$ is singular
\item $A$ is diagonalizable
\end{enumerate}

\textbf{Explanation:}
The characteristic polynomial has roots $\lambda = 0$ (algebraic
multiplicity~$2$), $\lambda = 3$, and $\lambda = -1$. Since $0$ is an
eigenvalue, there exists a nonzero $\mathbf{v}$ with
$A\mathbf{v} = \mathbf{0}$, so $\ker(A) \neq \{\mathbf{0}\}$.
Equivalently, $\det(A) = 0 \cdot 0 \cdot 3 \cdot (-1) = 0$, confirming
$A$ is singular.

\textbf{Trick answer:}
Since $\det(A) = 0$, $A^{-1}$ does \emph{not} exist --- the opposite of
what this choice claims.

\textbf{Trick answer:}
$\operatorname{tr}(A) = 0 + 0 + 3 + (-1) = 2 \neq 0$. A student who
confuses trace (sum of eigenvalues) with determinant (product) might
select this.

\textbf{Trick answer:}
$e^A$ is \emph{always} nonsingular, for any matrix $A$:
$\det(e^A) = e^{\operatorname{tr}(A)} = e^2 \neq 0$. This is a common
misconception: students think singularity of $A$ implies singularity
of $e^A$.

\textbf{Trick answer:}
The repeated eigenvalue $\lambda = 0$ has algebraic multiplicity~$2$,
but the geometric multiplicity could be $1$ or $2$. If geometric
multiplicity is~$1$, $A$ is not diagonalizable. This cannot be determined
from the characteristic polynomial alone.

\divider

%----------------------------------------------------------
\subsection*{Problem: Matrix powers via eigendecomposition}
%----------------------------------------------------------

A diagonalizable $2 \times 2$ matrix $A$ has eigenvalues $\lambda_1 = 2$
and $\lambda_2 = -1$ with corresponding eigenvectors $\mathbf{v}_1$
and $\mathbf{v}_2$.
Given $\mathbf{x}_0 = 3\mathbf{v}_1 + 5\mathbf{v}_2$,
what is $A^{100}\mathbf{x}_0$?

\begin{enumerate}[(A)]
\item $(3 \cdot 2^{100} + 5)\,\mathbf{v}_1$
\item $2^{100}(3\,\mathbf{v}_1 + 5\,\mathbf{v}_2)$
\item $3 \cdot 2^{100}\,\mathbf{v}_1 - 5\,\mathbf{v}_2$
\item $3 \cdot 2^{100}\,\mathbf{v}_1$
\item \textcolor{blue}{\textbf{None of the above}}
\end{enumerate}

\textbf{Explanation:}
When $\mathbf{x}_0$ is expanded in the eigenbasis, $A^k$ acts on each
component independently:
$$A^{100}\mathbf{x}_0 = 3\lambda_1^{100}\,\mathbf{v}_1
  + 5\lambda_2^{100}\,\mathbf{v}_2
  = 3 \cdot 2^{100}\,\mathbf{v}_1 + 5 \cdot (-1)^{100}\,\mathbf{v}_2$$
Since $(-1)^{100} = +1$ (even power), this simplifies to:
$$A^{100}\mathbf{x}_0 = 3 \cdot 2^{100}\,\mathbf{v}_1
  + 5\,\mathbf{v}_2$$
This expression does not appear among the other four choices, so the
answer is ``None of the above.'' The key ideas: (1)~each eigenvector
component is scaled by its own eigenvalue power, (2)~$(-1)^{100} = +1$,
and (3)~the two components evolve independently.

\textbf{Partial credit:}
The choice $3 \cdot 2^{100}\,\mathbf{v}_1 - 5\,\mathbf{v}_2$ correctly
applies the eigendecomposition formula and correctly recognizes that
$(-1)^{100}$ has magnitude~$1$. The only error is the sign:
$(-1)^{100} = +1$, not $-1$. This is a minor arithmetic slip reflecting
solid conceptual understanding with only the even/odd distinction missed.

\textbf{Trick answer:}
The choice $2^{100}(3\mathbf{v}_1 + 5\mathbf{v}_2)$ treats the dominant
eigenvalue $2^{100}$ as a uniform scaling factor applied to the entire
vector. This misses the fundamental point that each eigenvector component
is scaled by its \emph{own} eigenvalue power --- different directions are
stretched differently.

\textbf{Trick answer:}
The choice $3 \cdot 2^{100}\,\mathbf{v}_1$ assumes the $\mathbf{v}_2$
component vanishes because $(-1)^{100}$ is ``small'' or ``decays.'' In
fact, $|-1|^{100} = 1$: the magnitude is perfectly preserved. This
confuses $|\lambda| < 1$ (which causes decay) with $\lambda < 0$ (which
causes sign alternation but no decay).

\textbf{Trick answer:}
The choice $(3 \cdot 2^{100} + 5)\mathbf{v}_1$ collapses both eigenvector
components into a single direction, as if the large factor $2^{100}$
absorbs the $\mathbf{v}_2$ component. The components along different
eigenvectors are independent and never ``merge.''

\divider

%----------------------------------------------------------
\subsection*{Problem: Eigenvalues from trace and determinant}
%----------------------------------------------------------

A $3 \times 3$ real matrix $A$ satisfies $\operatorname{tr}(A) = 7$,
$\det(A) = 0$, and is known to have an eigenvalue $\lambda_1 = 3$.
Which of the following must be the remaining eigenvalues?

\begin{enumerate}[(A)]
\item \textcolor{blue}{\textbf{$\lambda_2 = 4,\; \lambda_3 = 0$}}
\item $\lambda_2 = 4,\; \lambda_3 = 3$
\item $\lambda_2 = 2+i,\; \lambda_3 = 2-i$
\item $\lambda_2 = 0,\; \lambda_3 = 0$
\item $\lambda_2 = 7,\; \lambda_3 = 0$
\end{enumerate}

\textbf{Explanation:}
Two constraints must be satisfied simultaneously:
(1)~$\operatorname{tr}(A) = \lambda_1 + \lambda_2 + \lambda_3 = 7$,
so $\lambda_2 + \lambda_3 = 4$;
(2)~$\det(A) = \lambda_1 \cdot \lambda_2 \cdot \lambda_3 = 0$, so at
least one of $\lambda_2, \lambda_3$ must be zero.
Additionally, complex eigenvalues of a real matrix come in conjugate pairs.
Checking: $\lambda_2 = 4$ and $\lambda_3 = 0$ satisfies both $4 + 0 = 4$
and $3 \cdot 4 \cdot 0 = 0$.

\textbf{Partial credit:}
The conjugate pair $2 \pm i$ satisfies the trace condition
($(2+i) + (2-i) = 4$) and the conjugate-pair requirement, but fails the
determinant: $3(2+i)(2-i) = 3 \cdot 5 = 15 \neq 0$. This student
correctly checked two out of three constraints.

\textbf{Trick answer:}
$\lambda_2 = 4, \lambda_3 = 3$ gives $3 + 4 + 3 = 10 \neq 7$ --- fails
the trace condition.

\textbf{Trick answer:}
$\lambda_2 = \lambda_3 = 0$ gives $0 + 0 = 0 \neq 4$ --- fails the
trace condition. Tempting for students who over-focus on $\det = 0$.

\textbf{Trick answer:}
$\lambda_2 = 7, \lambda_3 = 0$ gives $3 + 7 + 0 = 10 \neq 7$ --- fails
the trace condition. Appealing if one focuses only on
having a zero eigenvalue.

\divider

\newpage

%%%%%%%%%%%%%%%%%%%%%%%%%%%%%%%%%%%%%%%%%%%%%%%%%%%%%%%%%%%
\section*{WEEK 8: EIGENVALUE COMPLEXITIES}
%%%%%%%%%%%%%%%%%%%%%%%%%%%%%%%%%%%%%%%%%%%%%%%%%%%%%%%%%%%

%----------------------------------------------------------
\subsection*{Problem: Counting blocks in a real Jordan form}
%----------------------------------------------------------

The $6 \times 6$ matrix $\mathcal{J}$ below is in real Jordan canonical form:
$$\mathcal{J} = \begin{bmatrix}
3 & 1 & 0 & 0 & 0 & 0 \\
0 & 3 & 0 & 0 & 0 & 0 \\
0 & 0 & 3 & 0 & 0 & 0 \\
0 & 0 & 0 & 0 & 0 & 0 \\
0 & 0 & 0 & 0 & -1 & -2 \\
0 & 0 & 0 & 0 & 2 & -1
\end{bmatrix}$$
How many Jordan blocks does $\mathcal{J}$ have?

\begin{enumerate}[(A)]
\item $6$
\item $3$
\item \textcolor{blue}{\textbf{$4$}}
\item $5$
\item This matrix is not in Jordan canonical form
\end{enumerate}

\textbf{Explanation:}
The four blocks are:
\begin{enumerate}
\item A $2 \times 2$ Jordan block $\begin{bmatrix} 3 & 1 \\ 0 & 3 \end{bmatrix}$
      for $\lambda = 3$ (the superdiagonal~$1$ chains two vectors).
\item A $1 \times 1$ block $[3]$ for $\lambda = 3$.
\item A $1 \times 1$ block $[0]$ for $\lambda = 0$.
\item A $2 \times 2$ rotation block
      $\begin{bmatrix} -1 & -2 \\ 2 & -1 \end{bmatrix}$ encoding
      the complex conjugate pair $\lambda = -1 \pm 2i$.
\end{enumerate}
Key subtleties: (i)~the eigenvalue $\lambda = 3$ produces \emph{two}
separate blocks (a $2 \times 2$ chain and a $1 \times 1$), not one
combined block; (ii)~the $2 \times 2$ rotation block is a \emph{single}
block in real Jordan form, not two separate $1 \times 1$ blocks.

\textbf{Trick answer:}
The answer $3$ counts one block per distinct eigenvalue: one for
$\lambda = 3$, one for $\lambda = 0$, one for the complex pair. This
conflates ``same eigenvalue'' with ``same block.'' A single eigenvalue
can produce multiple blocks.

\textbf{Trick answer:}
The answer $5$ splits the $2 \times 2$ rotation block into two
$1 \times 1$ blocks by reading $-1$ and $-1$ off its diagonal. This
targets the student who does not recognize the rotation block as a
single unit encoding a complex conjugate pair.

\textbf{Trick answer:}
The answer $6$ treats every diagonal entry as its own block, ignoring all
block structure entirely.

\textbf{Trick answer:}
The claim that $\mathcal{J}$ is ``not in Jordan form'' targets a student
who only knows \emph{complex} Jordan form (where all blocks are upper
triangular) and does not recognize that \emph{real} Jordan form uses
$2 \times 2$ rotation blocks for complex conjugate pairs.

\divider

%----------------------------------------------------------
\subsection*{Problem: Number of Jordan blocks from eigenspace dimension}
%----------------------------------------------------------

A $4 \times 4$ matrix $A$ has only one eigenvalue, $\lambda = 3$, with
algebraic multiplicity~$4$. The eigenspace $\ker(A - 3I)$ has dimension~$2$.
How many Jordan blocks appear in the Jordan form of $A$, and why?

\begin{enumerate}[(A)]
\item Two blocks: specifically one $3 \times 3$ block and one $1 \times 1$
      block
\item Four $1 \times 1$ blocks, because the algebraic multiplicity is $4$
\item \textcolor{blue}{\textbf{Two blocks, because the number of Jordan
      blocks for an eigenvalue equals the dimension of its eigenspace}}
\item One $4 \times 4$ block, because there is only one eigenvalue
\item Cannot be determined without knowing $\dim(\ker(A - 3I)^k)$
      for $k \geq 2$
\end{enumerate}

\textbf{Explanation:}
A fundamental fact: the number of Jordan blocks for a given eigenvalue
equals the \emph{geometric multiplicity} --- the dimension of the
eigenspace. Each Jordan block begins with an eigenvector, so the
number of independent eigenvectors is exactly the number of blocks.
Here, $\dim(\ker(A - 3I)) = 2$, so there are exactly $2$ Jordan blocks
for $\lambda = 3$. The blocks must account for all $4$ dimensions;
they could be two $2 \times 2$ blocks or one $3 \times 3$ plus one
$1 \times 1$ --- but the \emph{count} is determined.

\textbf{Partial credit:}
The choice specifying one $3 \times 3$ and one $1 \times 1$ block
correctly identifies two blocks (the key insight). The error is
overcommitting to a specific partition when two $2 \times 2$ blocks is
equally consistent with the given data. Good understanding, but goes
one step further than the information allows.

\textbf{Trick answer:}
Four $1 \times 1$ blocks would mean $A$ is diagonalizable. But
geometric multiplicity~$2 <$ algebraic multiplicity~$4$ means $A$ is
\emph{not} diagonalizable.

\textbf{Trick answer:}
One $4 \times 4$ block confuses ``one eigenvalue'' with ``one block.''
Having a single eigenvalue constrains all blocks to share that eigenvalue,
but the block count depends on the eigenspace dimension, not the number of
distinct eigenvalues.

\textbf{Trick answer:}
The block \emph{count} is determined by the eigenspace dimension alone.
While generalized eigenspaces are needed for block \emph{sizes}, the
question only asks about the number of blocks.

\divider

%----------------------------------------------------------
\subsection*{Problem: Complex eigenvalues --- roles of real and imaginary parts}
%----------------------------------------------------------

A $2 \times 2$ real matrix $A$ has complex eigenvalues
$\lambda = \alpha \pm i\beta$ with $\alpha < 0$ and $\beta > 0$.
Solutions to $\dfrac{d\mathbf{x}}{dt} = A\mathbf{x}$ oscillate with
decreasing amplitude.
Now consider $A'$ with eigenvalues $\lambda' = \alpha \pm 2i\beta$.
Compared to the original, solutions of
$\dfrac{d\mathbf{x}}{dt} = A'\mathbf{x}$ will:

\begin{enumerate}[(A)]
\item Oscillate at twice the frequency and also decay faster, since the
      eigenvalues are farther from the origin in the complex plane
\item Have identical behavior, since only the real part affects a real ODE
\item Oscillate at the same frequency but decay faster, since the imaginary
      part is larger
\item Oscillate at twice the frequency and decay at half the rate, since
      energy is transferred from decay into oscillation
\item \textcolor{blue}{\textbf{Oscillate at twice the frequency with the
      same rate of amplitude decay}}
\end{enumerate}

\textbf{Explanation:}
For eigenvalues $\alpha \pm i\beta$, the real basis solutions have the form
$e^{\alpha t}\cos(\beta t)$ and $e^{\alpha t}\sin(\beta t)$.
The real part $\alpha$ controls exponential decay; the imaginary part $\beta$
controls oscillation frequency. These roles are \emph{completely independent}.
Changing $\beta$ to $2\beta$ doubles the oscillation frequency; $\alpha$ is
unchanged, so the decay rate is identical.

\textbf{Partial credit:}
The choice claiming ``twice the frequency and faster decay'' correctly
identifies the frequency doubling but incorrectly infers faster decay from
$|\lambda'| > |\lambda|$. This reflects a confusion between eigenvalue
\emph{magnitude} and decay rate --- only the real part governs decay.

\textbf{Trick answer:}
The choice ``same frequency, faster decay'' reverses which part does what.

\textbf{Trick answer:}
The ``energy transfer'' choice invokes a plausible-sounding but fictitious
conservation principle. There is no energy trade-off between oscillation
and decay; $\alpha$ and $\beta$ act independently.

\textbf{Trick answer:}
The choice ``identical behavior'' confuses the fact that we work with
\emph{real} solutions with the claim that the imaginary part is irrelevant.
The imaginary part determines the oscillation frequency even though the
solution is real-valued.

\divider

%----------------------------------------------------------
\subsection*{Problem: Higher-order ODE and basis solutions}
%----------------------------------------------------------

A fourth-order linear homogeneous ODE has characteristic polynomial
$$p(\lambda) = (\lambda - 1)(\lambda + 2)^2(\lambda + 3) = 0$$
Which of the following is the correct general solution?

\begin{enumerate}[(A)]
\item $x(t) = c_1 e^{t} + c_2 e^{-2t} + c_3 e^{-3t}$
\item $x(t) = c_1 e^{t} + c_2 e^{-2t} + c_3 e^{-2t}\cos(t) + c_4 e^{-3t}$
\item \textcolor{blue}{\textbf{$x(t) = c_1 e^{t} + c_2 e^{-2t}
      + c_3 te^{-2t} + c_4 e^{-3t}$}}
\item $x(t) = c_1 e^{t} + c_2 e^{-2t} + c_3 t^2 e^{-2t} + c_4 e^{-3t}$
\item $x(t) = c_1 e^{t} + c_2 e^{-2t} + c_3 te^{-2t} + c_4 te^{-3t}$
\end{enumerate}

\textbf{Explanation:}
The roots are $\lambda = 1$ (simple), $\lambda = -2$ (multiplicity~$2$),
and $\lambda = -3$ (simple). The rules for basis solutions:
each simple root $\lambda$ contributes $e^{\lambda t}$;
a repeated root $\lambda$ with multiplicity~$k$ contributes
$e^{\lambda t},\, te^{\lambda t},\, \ldots,\, t^{k-1}e^{\lambda t}$.
Here $\lambda = -2$ appears twice, producing $e^{-2t}$ and $te^{-2t}$.
Combined with $e^{t}$ and $e^{-3t}$, we get four basis solutions matching
the fourth-order equation.

\textbf{Trick answer:}
The choice with only $3$ terms ignores the multiplicity of $\lambda = -2$
entirely, producing too few basis solutions for a fourth-order equation.

\textbf{Trick answer:}
The choice involving $\cos(t)$ confuses repeated \emph{real} eigenvalues
with \emph{complex} eigenvalues. Trigonometric terms arise from imaginary
parts; $\lambda = -2$ is purely real.

\textbf{Trick answer:}
The choice with $t^2 e^{-2t}$ uses the wrong power of $t$: for
multiplicity~$k$, the prefactors run $t^0, t^1, \ldots, t^{k-1}$, so
multiplicity~$2$ gives $e^{-2t}$ and $te^{-2t}$, not $t^2 e^{-2t}$.

\textbf{Trick answer:}
The choice with both $te^{-2t}$ and $te^{-3t}$ correctly handles the
repeated root but applies a spurious polynomial prefactor to the
\emph{simple} root $\lambda = -3$. Polynomial prefactors arise only for
repeated eigenvalues.

\divider

%----------------------------------------------------------
\subsection*{Problem: QR algorithm and eigenvalue preservation}
%----------------------------------------------------------

The QR algorithm computes $A_{k+1} = R_k Q_k$ from the QR factorization
$A_k = Q_k R_k$.
A key property is that every iterate $A_{k+1}$ has the same eigenvalues
as $A_k$. Why is this true?

\begin{enumerate}[(A)]
\item Because $Q_k$ is orthogonal, and orthogonal transformations
      preserve all matrix properties including eigenvalues
\item Because the QR factorization preserves eigenvalues:
      $Q_k$ and $R_k$ each have the same eigenvalues as $A_k$
\item \textcolor{blue}{\textbf{Because $A_{k+1} = Q_k^{-1} A_k Q_k$,
      which is a similarity transformation}}
\item Because the diagonal entries of $R_k$ are the eigenvalues of
      $A_k$, and these entries are reused to construct $A_{k+1}$
\item Because $R_k Q_k$ and $Q_k R_k$ are transposes of each other,
      and a matrix and its transpose share eigenvalues
\end{enumerate}

\textbf{Explanation:}
The crucial identity:
$$A_{k+1} = R_k Q_k = (Q_k^{-1} Q_k) R_k Q_k
  = Q_k^{-1}(Q_k R_k) Q_k = Q_k^{-1} A_k Q_k$$
So $A_{k+1}$ is obtained from $A_k$ by a similarity transformation.
Similar matrices have the same characteristic polynomial and therefore
the same eigenvalues. This is the foundational reason the QR algorithm
works: it iteratively transforms $A$ toward a simpler form while
preserving eigenvalues at every step.

\textbf{Trick answer:}
``Orthogonal transformations preserve all matrix properties'' is far too
broad. Eigenvalue preservation comes specifically from the \emph{similarity}
$Q_k^{-1}A_kQ_k$; any invertible $P$ gives a similarity $P^{-1}A_kP$
preserving eigenvalues. Orthogonality helps with numerical stability but
is not the reason eigenvalues are preserved.

\textbf{Trick answer:}
The QR factorization does \emph{not} individually preserve eigenvalues.
The eigenvalues of $Q_k$ all have modulus~$1$; the eigenvalues of $R_k$
are its diagonal entries. Neither matches $A_k$'s eigenvalues in general.

\textbf{Trick answer:}
The diagonal entries of $R_k$ come from norms in the Gram--Schmidt process,
not from the eigenvalues of $A_k$. In the Schur form (obtained by
similarity), diagonal entries \emph{are} eigenvalues --- but $R$ from
QR factorization is a different object.

\textbf{Trick answer:}
$R_kQ_k$ and $Q_kR_k$ are generally \emph{not} transposes of each other.
While it is true that a matrix and its transpose share eigenvalues, that
fact is irrelevant here because the transpose relationship does not hold.

\divider

%----------------------------------------------------------
\subsection*{Problem: Sufficient condition for diagonalizability}
%----------------------------------------------------------

A $3 \times 3$ matrix $A$ has eigenvalues $\lambda_1 = 4$,
$\lambda_2 = 4$, $\lambda_3 = 1$.
Which additional piece of information would guarantee that $A$ is
diagonalizable?

\begin{enumerate}[(A)]
\item $\ker(A - 4I)$ is one-dimensional
\item $A$ is upper triangular
\item The columns of $A$ are linearly independent
\item \textcolor{blue}{\textbf{The eigenspace for $\lambda = 4$ is
      two-dimensional}}
\item None of the above: a repeated eigenvalue means $A$ cannot
      be diagonalized
\end{enumerate}

\textbf{Explanation:}
The eigenvalue $\lambda = 4$ has algebraic multiplicity~$2$, so
diagonalizability hinges on whether its geometric multiplicity also
equals~$2$. If the eigenspace is two-dimensional, geometric $=$
algebraic for this eigenvalue. The eigenvalue $\lambda = 1$ has
algebraic multiplicity~$1$, so its geometric multiplicity is
automatically~$1$. With geometric $=$ algebraic for every eigenvalue,
$A$ is diagonalizable.

\textbf{Trick answer:}
A one-dimensional eigenspace for $\lambda = 4$ means geometric
multiplicity~$1 < 2 =$ algebraic multiplicity, which guarantees
$A$ is \emph{not} diagonalizable. The student has the right framework
but the wrong dimension threshold.

\textbf{Trick answer:}
Upper triangular matrices display their eigenvalues on the diagonal,
but this gives no information about eigenspace dimensions. A Jordan block
like $\begin{bmatrix} 4 & 1 \\ 0 & 4 \end{bmatrix}$ is upper triangular
with repeated eigenvalue~$4$ yet is not diagonalizable.

\textbf{Trick answer:}
Since no eigenvalue is zero,
$\det(A) = 4 \cdot 4 \cdot 1 = 16 \neq 0$, so $A$ is already
invertible (columns linearly independent) from the given information
alone. This tells us nothing new, and invertibility is unrelated to
diagonalizability.

\textbf{Trick answer:}
Repeated eigenvalues do \emph{not} prevent diagonalizability; they
merely create a condition that must be checked. For instance,
$4I_{3 \times 3}$ has eigenvalue~$4$ with multiplicity~$3$ and is
trivially diagonalizable (it is already diagonal).

\divider

\newpage

%%%%%%%%%%%%%%%%%%%%%%%%%%%%%%%%%%%%%%%%%%%%%%%%%%%%%%%%%%%
\section*{WEEK 9: LINEAR ITERATIVE SYSTEMS}
%%%%%%%%%%%%%%%%%%%%%%%%%%%%%%%%%%%%%%%%%%%%%%%%%%%%%%%%%%%

%----------------------------------------------------------
\subsection*{Problem: Spectral radius and convergence to zero}
%----------------------------------------------------------

A $3 \times 3$ matrix $A$ has eigenvalues $\lambda_1 = 0.7$,
$\lambda_2 = -0.9$, and $\lambda_3 = 0.2$.
Consider the iteration $\mathbf{x}_{k+1} = A\mathbf{x}_k$.
What is the long-term behavior of $\mathbf{x}_k$?

\begin{enumerate}[(A)]
\item $\mathbf{x}_k$ may converge or diverge depending on the initial
      condition
\item $\mathbf{x}_k$ oscillates indefinitely without settling to a limit
\item $\mathbf{x}_k$ converges to a dominant eigenvector
\item \textcolor{blue}{\textbf{$\mathbf{x}_k$ converges to $\mathbf{0}$
      for every initial condition}}
\item $\mathbf{x}_k$ diverges due to oscillations at each step
\end{enumerate}

\textbf{Explanation:}
$A^k \to 0$ if and only if every eigenvalue satisfies $|\lambda_i| < 1$,
i.e., the spectral radius $\rho(A) < 1$. Here $|\lambda_1| = 0.7$,
$|\lambda_2| = 0.9$, $|\lambda_3| = 0.2$, so $\rho(A) = 0.9 < 1$.
The component $(-0.9)^k$ does alternate in sign, but its magnitude
$0.9^k \to 0$: the oscillations are damped and die out. Every component
decays, so $\mathbf{x}_k \to \mathbf{0}$ for \emph{all} initial conditions.

\textbf{Partial credit:}
The ``oscillates indefinitely'' choice correctly identifies the
sign-alternating behavior of $(-0.9)^k$, but concludes the oscillation
persists at constant amplitude. In fact, $0.9^k \to 0$, so the
oscillations decay. The student understands the qualitative effect of a
negative eigenvalue but fails to connect $|\lambda| < 1$ with decay.

\textbf{Trick answer:}
``Diverges due to oscillations'' confuses the \emph{sign} of $\lambda_2$
with growth. A negative eigenvalue with $|\lambda| < 1$ produces
\emph{decaying} sign-alternating behavior, not growing oscillations.

\textbf{Trick answer:}
``Converges to a dominant eigenvector'' makes two compounded errors:
(i)~identifies $\lambda_1 = 0.7$ as ``the largest'' by value instead of
by magnitude ($|\lambda_2| = 0.9 > 0.7$), and (ii)~claims convergence to
a nonzero limit, which would require $|\lambda| = 1$, not $< 1$.

\textbf{Trick answer:}
``Depends on the initial condition'' arises when some eigenvalues have
$|\lambda| < 1$ and others have $|\lambda| > 1$. Here \emph{all}
$|\lambda_i| < 1$, so convergence to $\mathbf{0}$ is universal.

\divider

%----------------------------------------------------------
\subsection*{Problem: Graph Laplacian and the zero eigenvalue}
%----------------------------------------------------------

Let $G$ be an undirected graph on $n$ vertices with adjacency matrix $A$
and degree matrix $D = \operatorname{diag}(d_1, \ldots, d_n)$. The graph
Laplacian is $L = D - A$.
Which best explains why $\lambda = 0$ is always an eigenvalue of $L$?

\begin{enumerate}[(A)]
\item \textcolor{blue}{\textbf{Every row of $L$ sums to zero, so the
      all-ones vector $\mathbf{1}$ is in $\ker(L)$}}
\item $L = D - A$ is a difference of two matrices, and the subtraction
      causes cancellation that forces a zero eigenvalue
\item The diagonal entries of $L$ are non-negative and the off-diagonal
      entries are non-positive, which forces at least one eigenvalue to
      be zero
\item $L$ is symmetric, and symmetric matrices always have at least one
      zero eigenvalue
\item None of the above
\end{enumerate}

\textbf{Explanation:}
Each row of $L$ sums to zero because
$\sum_j L_{ij} = d_i - \sum_{j \neq i} A_{ij} = d_i - d_i = 0$.
Therefore $L\mathbf{1} = \mathbf{0}$, making $\mathbf{1}$ an eigenvector
with eigenvalue~$0$. This works for any undirected graph.

\textbf{Trick answer:}
Symmetric matrices do \emph{not} necessarily have a zero eigenvalue.
The identity matrix $I$ is symmetric with all eigenvalues equal to~$1$.
Symmetry guarantees real eigenvalues and orthogonal eigenvectors, but
nothing about zero being among them.

\textbf{Trick answer:}
That $L$ is a ``difference'' does not explain the zero eigenvalue.
Many differences are invertible --- e.g., $3I - I = 2I$ has no zero
eigenvalue. The zero eigenvalue arises from the specific algebraic
relationship between $D$ and $A$.

\textbf{Trick answer:}
The sign pattern of $L$ does not force a zero eigenvalue.
$\begin{bmatrix} 2 & -1 \\ -1 & 3 \end{bmatrix}$ has the same sign
pattern but is invertible. The sign pattern describes M-matrices, not
all of which are singular.

\divider

%----------------------------------------------------------
\subsection*{Problem: Spectral Theorem vs.\ diagonalizability}
%----------------------------------------------------------

A real $n \times n$ matrix $A$ is known to be symmetric.
Which of the following is a consequence of the Spectral Theorem that does
\textbf{not} follow from diagonalizability alone?

\begin{enumerate}[(A)]
\item $\operatorname{tr}(A)$ equals the sum of the eigenvalues of $A$
\item $\det(A)$ equals the product of the eigenvalues of $A$
\item $A$ can be written as $V\Lambda V^{-1}$ for some invertible $V$
      and diagonal $\Lambda$
\item $A$ has $n$ eigenvalues with algebraic and geometric multiplicity one
\item \textcolor{blue}{\textbf{The eigenvectors of $A$ can be chosen to
      form an orthonormal basis for $\mathbb{R}^n$}}
\end{enumerate}

\textbf{Explanation:}
The Spectral Theorem states $A = Q\Lambda Q^T$ with $Q$ orthogonal and
$\Lambda$ real diagonal. The key upgrade over mere diagonalizability is
the existence of an \emph{orthonormal} eigenbasis. Diagonalizability
gives a basis of eigenvectors; the Spectral Theorem upgrades this to an
orthonormal one.

Among the other choices:
trace $=$ sum of eigenvalues and det $=$ product of eigenvalues hold for
\emph{all} square matrices; $A = V\Lambda V^{-1}$ is the
\emph{definition} of diagonalizability; and the claim about multiplicities
is false (symmetric matrices can have repeated eigenvalues).

\textbf{Trick answer:}
$A = V\Lambda V^{-1}$ \emph{is} diagonalizability --- precisely the
property the question says is already assumed. A student who selects this
fails to distinguish diagonalizability from the additional structure
symmetry provides.

\textbf{Trick answer:}
The claim about $n$ eigenvalues each with algebraic and geometric
multiplicity one is false --- symmetric matrices can have repeated
eigenvalues. The Spectral Theorem guarantees a full orthonormal
eigenbasis even with repetition.

\textbf{Trick answer:}
The trace and determinant facts hold universally for all square matrices
and have nothing to do with diagonalizability or symmetry.

\divider

%----------------------------------------------------------
\subsection*{Problem: Positive stochastic matrix and stationary distribution}
%----------------------------------------------------------

Consider the $3 \times 3$ positive stochastic matrix $P$ with stationary
distribution $\boldsymbol{\pi} = (0.2,\, 0.4,\, 0.4)^T$.
If $\mathbf{x}_0 = (1,\, 0,\, 0)^T$, what is
$\displaystyle\lim_{k \to \infty} P^k \mathbf{x}_0$?

\begin{enumerate}[(A)]
\item \textcolor{blue}{$\boldsymbol{\begin{pmatrix} 0.2 \\ 0.4 \\ 0.4 \end{pmatrix}}$}
\item The limit depends on $\mathbf{x}_0$; knowing
      $\boldsymbol{\pi}$ alone is not enough
\item $\begin{pmatrix} 1/3 \\ 1/3 \\ 1/3 \end{pmatrix}$
\item $\begin{pmatrix} 0.4 \\ 0.4 \\ 0.2 \end{pmatrix}$
\item The limit does not necessarily exist
\end{enumerate}

\textbf{Explanation:}
By Perron--Frobenius, a positive stochastic matrix has a unique stationary
distribution $\boldsymbol{\pi}$, and all other eigenvalues satisfy
$|\lambda_i| < 1$. The subdominant components of $\mathbf{x}_0$ decay,
leaving $P^k\mathbf{x}_0 \to \boldsymbol{\pi}$ for \emph{any} initial
probability distribution. The initial state affects transient behavior
but not the limit.

\textbf{Trick answer:}
The choice $(0.4,\, 0.4,\, 0.2)^T$ is $P\mathbf{x}_0$ --- the first
column of $P$, i.e., the distribution after \emph{one} step, not
infinitely many. Particularly deceptive because the entries $0.4, 0.4, 0.2$
are the same three numbers as $\boldsymbol{\pi} = (0.2, 0.4, 0.4)^T$
rearranged.

\textbf{Trick answer:}
The uniform distribution $(1/3,\, 1/3,\, 1/3)^T$ confuses ``equilibrium''
with ``uniform.'' The stationary distribution reflects the structure of
$P$; uniform is stationary only for doubly stochastic matrices.

\textbf{Trick answer:}
``The limit depends on $\mathbf{x}_0$'' sounds sophisticated but is wrong
for a positive stochastic matrix. Perron--Frobenius guarantees a
\emph{unique} stationary distribution that attracts every probability
distribution.

\textbf{Trick answer:}
``The limit does not exist'' reverses the logic of stationarity. The
convergence theorem says much more than $\boldsymbol{\pi}$ being a fixed
point: it says $\boldsymbol{\pi}$ is a global attractor.

\divider

%----------------------------------------------------------
\subsection*{Problem: Stochastic matrix eigenvalue properties}
%----------------------------------------------------------

A $3 \times 3$ matrix $P$ has non-negative entries and every column
sums to~$1$.
Which of the following \textbf{must} be true?

\begin{enumerate}[(A)]
\item All eigenvalues of $P$ are real
\item The columns of $P$ are linearly independent
\item Every eigenvalue of $P$ satisfies $0 \leq \lambda \leq 1$
\item $\det(P) = 1$
\item \textcolor{blue}{\textbf{None of the above}}
\end{enumerate}

\textbf{Explanation:}
Each of (A)--(D) fails for some valid stochastic matrix. The properties
that \emph{do} always hold are: $\lambda = 1$ is always an eigenvalue,
and every eigenvalue satisfies $|\lambda| \leq 1$ (a bound on
\emph{magnitude}). None of the four listed statements captures these
correctly.

\textbf{Partial credit:}
The choice $0 \leq \lambda \leq 1$ is the most instructive wrong answer.
The correct bound is $|\lambda| \leq 1$ --- a constraint on magnitude.
This allows negative real eigenvalues and complex eigenvalues, as long as
the modulus does not exceed~$1$. The student recalls the spectral radius
bound but misinterprets it as a pointwise inequality on the real line.

\textbf{Trick answer:}
Stochastic matrices can have complex eigenvalues. The cyclic permutation
matrix $\begin{bmatrix} 0&0&1\\1&0&0\\0&1&0 \end{bmatrix}$ is stochastic
with eigenvalues $1,\, e^{2\pi i/3},\, e^{-2\pi i/3}$ --- non-real.

\textbf{Trick answer:}
$\det(P)$ equals the product of eigenvalues, and there is no reason this
product must be~$1$. A stochastic matrix with eigenvalues $1, 0.5, -0.3$
has $\det(P) = -0.15$.

\textbf{Trick answer:}
Stochastic matrices can have linearly dependent columns (e.g., two
identical columns), making them singular. A student may assume probability
transition matrices should always be invertible.

\divider

%----------------------------------------------------------
\subsection*{Problem: Perron--Frobenius theorem guarantees}
%----------------------------------------------------------

A positive stochastic matrix $P$ (all entries strictly positive, columns
sum to~$1$) has eigenvalues $\lambda_i$.
The Perron--Frobenius theorem guarantees several properties of $P$.
Which of the following is \textbf{not} among them?

\begin{enumerate}[(A)]
\item Every other eigenvalue satisfies $|\lambda_i| < 1$
\item The eigenvalue $\lambda_1 = 1$ is simple (algebraic multiplicity one)
\item The eigenvalue $\lambda_1 = 1$ has a corresponding eigenvector with
      all entries strictly positive
\item For any initial probability distribution $\mathbf{x}_0$, the iterates
      $P^k\mathbf{x}_0$ converge to the unique stationary distribution
\item \textcolor{blue}{\textbf{None of the above: all these properties are
      guaranteed}}
\end{enumerate}

\textbf{Explanation:}
All four listed properties are indeed guaranteed by Perron--Frobenius for
a \emph{positive} stochastic matrix:
(A)~the spectral gap $|\lambda_i| < 1$ for $i \geq 2$;
(B)~simplicity of the Perron eigenvalue;
(C)~positivity of the Perron eigenvector;
(D)~convergence to the unique stationary distribution from any starting
distribution.

\textbf{Trick answer:}
Each of the four properties targets a different aspect of Perron--Frobenius
that a student might doubt:
the spectral gap (perhaps confusing positive with merely non-negative
matrices), simplicity of the Perron eigenvalue (perhaps recalling weaker
results for reducible matrices), positivity of the eigenvector (perhaps
doubting that an eigenvector can have a sign constraint), or universal
convergence (perhaps thinking initial conditions affect the limit).
All are guaranteed for \emph{positive} matrices.

\divider

%----------------------------------------------------------
\subsection*{Problem: Orthonormal eigenbasis and dot-product coordinates}
%----------------------------------------------------------

Let $A$ be a real symmetric $n \times n$ matrix, and let $B$ be a general
(non-symmetric) diagonalizable $n \times n$ matrix. Both have $n$ linearly
independent eigenvectors. To express a vector $\mathbf{x} \in \mathbb{R}^n$
in the eigenbasis of each matrix, what key simplification does the symmetry
of $A$ provide?

\begin{enumerate}[(A)]
\item The eigenvalues of $A$ are guaranteed real, so the eigenbasis
      coordinates are real-valued
\item Both require computing $V^{-1}\mathbf{x}$; symmetry does not
      simplify this step
\item The diagonalization of $A$ is unique, whereas $B = V\Lambda V^{-1}$
      has infinitely many choices of $V$
\item \textcolor{blue}{\textbf{The eigenvectors of $A$ can always be
      chosen orthonormal, so the eigenbasis coordinates are
      $c_i = \mathbf{q}_i^T\mathbf{x}$ --- no matrix inversion is
      required}}
\item Symmetric matrices always have distinct eigenvalues, so there is
      exactly one eigenbasis up to scaling
\end{enumerate}

\textbf{Explanation:}
The Spectral Theorem guarantees an orthonormal eigenbasis
$\{\mathbf{q}_1, \ldots, \mathbf{q}_n\}$ with $Q^{-1} = Q^T$.
Eigenbasis coordinates become:
$c_i = \mathbf{q}_i^T\mathbf{x}$ --- a dot product, an $O(n)$ operation.
For a general diagonalizable matrix, coordinates require solving
$V^{-1}\mathbf{x}$, which is $O(n^2)$ or $O(n^3)$ and may be
ill-conditioned. The student must supply the chain:
symmetric $\Rightarrow$ Spectral Theorem $\Rightarrow$ orthonormal
eigenbasis $\Rightarrow$ $Q^{-1} = Q^T$ $\Rightarrow$ coordinates are
dot products.

\textbf{Partial credit:}
Real eigenvalues are indeed guaranteed for symmetric matrices, but this
does not address the computational simplification. A general real matrix
can also have all real eigenvalues and still require $V^{-1}$.

\textbf{Trick answer:}
The claim that symmetry does not simplify is precisely wrong --- this is
exactly what the Spectral Theorem simplifies. A student who does not know
the theorem may conclude both cases are ``just diagonalizable.''

\textbf{Trick answer:}
Diagonalization is not unique for either symmetric or general matrices ---
within each eigenspace, any basis may be chosen.

\textbf{Trick answer:}
Symmetric matrices can have repeated eigenvalues. The Spectral Theorem
guarantees a full orthonormal eigenbasis even with repetition.

\divider


\end{document}